\section{智能体的训练：合成数据与探索}

\subsection{当前实践：上下文学习}

当前训练 LLM 智能体的标准做法是：为每个新环境提供\textbf{人工编写的 few-shot 示范}作为上下文。

\begin{figure}[H]
\centering
\includegraphics[width=0.95\textwidth]{slides-images/slide-055.jpg}
\caption{当前实践的局限：需要为每种环境提供人工示范，面对数千种环境和无数种交互方式，这种方法不可扩展\protect\footnotemark}
\end{figure}
\footnotetext{Slides 第56页。}

\begin{warningbox}{人工示范不可扩展}
现实中有成千上万种不同的网站和应用，每种环境的交互方式都不同。为每种环境都准备人工示范既昂贵又不现实。核心问题变为：\textbf{智能体能否自主探索环境来生成高质量的合成训练数据？}
\end{warningbox}

\subsection{BAGEL：探索 + 模型生成数据}

BAGEL 方法巧妙地结合了前面在推理部分学到的\textbf{迭代自我训练}思想，将其应用到智能体的训练数据生成中。

\begin{figure}[H]
\centering
\includegraphics[width=0.95\textwidth]{slides-images/slide-060.jpg}
\caption{探索 + 模型生成数据：语言模型根据指令执行动作序列，通过过滤器判断轨迹质量，失败的轨迹由标签模型重新标注\protect\footnotemark}
\end{figure}
\footnotetext{Slides 第61页。}

BAGEL 的完整流程：

\begin{enumerate}
    \item \textbf{探索}：让语言模型\textbf{无目标地随机探索}环境，收集动作轨迹
    \item \textbf{标注}：用第二个语言模型为轨迹生成\textbf{自然语言描述}——推断``这条轨迹完成了什么任务''
    \item \textbf{条件生成}：用推断出的描述作为指令，让模型\textbf{有目标地}重新生成轨迹
    \item \textbf{过滤}：用过滤器判断（指令，轨迹）对的质量
    \item \textbf{重标注}：对于失败的轨迹，重新推断它实际完成了什么（而非原目标），赋予新标签
    \item \textbf{迭代}：重复步骤3--5，直到获得足够多的高质量（指令，轨迹）对
\end{enumerate}

\begin{figure}[H]
\centering
\includegraphics[width=0.95\textwidth]{slides-images/slide-065.jpg}
\caption{BAGEL 系统概览：三个阶段——(1) 探索环境收集轨迹；(2) 通过迭代重标注创建合成示范；(3) 推理时检索相关示范作为上下文\protect\footnotemark}
\end{figure}
\footnotetext{Slides 第66页。}

\begin{importantbox}{BAGEL 的核心创新：重标注}
当智能体未能完成目标任务时，BAGEL 不会丢弃这条轨迹，而是\textbf{重新标注}它——推断智能体实际完成了什么子任务。例如，智能体试图``订一张从旧金山到纽约的机票''但只完成了``将出发地设为 SFO、目的地设为纽约''，重标注后这条轨迹仍然是有价值的训练数据。这与推理部分的迭代自我训练异曲同工。
\end{importantbox}

获得合成训练数据后，有两种使用方式：

\begin{itemize}
    \item \textbf{上下文学习}：用合成示范替代人工示范作为 few-shot 示例
    \item \textbf{微调}：直接在合成数据上微调模型
\end{itemize}

实验表明，在 MiniWoB 上使用 BAGEL 可以获得 \textbf{13个百分点}的提升。

\subsection{本章小结}

通过探索环境和迭代重标注，智能体可以自主生成高质量的合成训练数据，减少对人工示范的依赖。这种思路与推理部分的自我训练一脉相承——核心都是让模型从自己的经验中学习和改进。

%% ========== Part 10: 多模态与挑战 ==========

